% !TEX root = ../main.tex
% Figure content only: inserted inside an outer figure environment in ch09
\begin{tikzpicture}[x=1cm,y=1cm]
  \tikzset{
    wavenode/.style={circle,draw=blue!70,thick,fill=blue!10,minimum size=1.05cm,font=\small},
    coup1/.style={<->,>=Stealth,thick,red!75},
    coup2/.style={<->,>=Stealth,thick,orange!85!black},
    loss/.style={->,>=Stealth,thick,dashed,red!60},
    gain/.style={->,>=Stealth,thick,green!50!black}
  }

  % A
  \begin{scope}[xshift=0cm,yshift=3.5cm]
    \node[font=\bfseries] at (0,2.2) {(a) 四波基底};
    \node[wavenode] (rx) at (-1.8,0.9) {$R_x$};
    \node[wavenode] (sx) at (1.8,0.9) {$S_x$};
    \node[wavenode] (ry) at (-1.8,-0.9) {$R_y$};
    \node[wavenode] (sy) at (1.8,-0.9) {$S_y$};
    \draw[coup1] (rx) -- node[above,font=\scriptsize] {$\kappa_1$} (sx);
    \draw[coup1] (ry) -- node[below,font=\scriptsize] {$\kappa_1$} (sy);
    \draw[coup2] (rx) -- node[left,font=\scriptsize] {$\kappa_2$} (ry);
    \draw[coup2] (sx) -- node[right,font=\scriptsize] {$\kappa_2$} (sy);
    \draw[coup2] (rx) -- (sy);
    \draw[coup2] (sx) -- (ry);
    \draw[loss] (rx) -- ++(0,0.9) node[above,font=\scriptsize] {$\gamma_{\rm rad}$};
    \draw[loss] (sx) -- ++(0,0.9);
    \draw[loss] (ry) -- ++(0,-0.9) node[below,font=\scriptsize] {$\gamma_{\rm rad}$};
    \draw[loss] (sy) -- ++(0,-0.9);
    \draw[gain] (-3.3,0.9) -- node[above,font=\scriptsize] {$g_{\rm mod}$} (rx);
    \draw[gain] (3.3,0.9) -- (sx);
    \draw[gain] (-3.3,-0.9) -- node[below,font=\scriptsize] {$g_{\rm mod}$} (ry);
    \draw[gain] (3.3,-0.9) -- (sy);
  \end{scope}

  % B
  \begin{scope}[xshift=-3.8cm,yshift=-1.2cm]
    \node[font=\bfseries] at (2.2,1.7) {(b) 矩阵形式};
    \node[align=center,font=\small] at (2.2,-0.1) {
      $\displaystyle
      \frac{\dd}{\dd t}
      \begin{bmatrix}
      R_x\\ S_x\\ R_y\\ S_y
      \end{bmatrix}
      =
      \left(\ii\bm{\Omega}-\bm{\Gamma}+\bm{G}\right)
      \begin{bmatrix}
      R_x\\ S_x\\ R_y\\ S_y
      \end{bmatrix}
      $
    };
    \node[draw=black,rounded corners,fill=gray!6,align=left,font=\footnotesize,text width=4.8cm] at (2.2,-2.1) {
      \textbf{读法：}
      $\bm{\Omega}$ 负责保守耦合，
      $\bm{\Gamma}$ 负责辐射/吸收损耗，
      $\bm{G}$ 负责模增益。
    };
  \end{scope}

  % C
  \begin{scope}[xshift=4.8cm,yshift=-1.2cm]
    \node[font=\bfseries] at (1.8,1.7) {(c) 信号流解释};
    \node[draw=blue!70,rounded corners,fill=blue!6,align=left,font=\footnotesize,text width=4.6cm] at (1.8,-0.5) {
      \textbf{$\kappa_1$：}反向散射，类似 1D DFB 中的前后向反馈。\\[3pt]
      \textbf{$\kappa_2$：}转角散射，把 $x$ 向和 $y$ 向波连接起来，是 2D PCSEL 的关键。\\[3pt]
      \textbf{$\gamma_{\rm rad}$：}通向光锥的辐射泄漏。\\[3pt]
      \textbf{$g_{\rm mod}$：}有源区对四波基底的等效增益注入。
    };
  \end{scope}
\end{tikzpicture}
